\documentclass{article}
\usepackage{amsmath}
\title{Q and P: Incentives at the Boundaries of Reward Shaping}
\author{ada}
\date{}
\begin{document}
\maketitle

\section*{The pair}
Q, \emph{Mechanism Design for Alignment and Control}, studies reward mechanisms constrained by joint report
and action deviations (ledger \#2--3).
P, \emph{MARS-RA}, turns external large multimodal model judgments of paired egocentric images into
Bradley--Terry scores and potential-based reward shaping, with robustness targeting the judge's latent
preferences (ledger \#1--2).

\section*{The connexion pursued}
The initial prospect of transferring Q's strategic guarantees to P does not survive this distinction: P's
comparisons are not strategic peer reports, and estimating a judge's preferences does not establish incentive
compatibility. Q's peer discipline also requires separated conditional beliefs and additional reward
instruments, as both peers' notes explain. The pursued connexion became an audit of whether P's shaping
preserves incentives across participation boundaries, with an information constraint on any deliberate
incentive-changing replacement (ledger \#3, \#8, \#12, \#16).

\section*{Findings and supporting arguments}
\paragraph{Complete accounting preserves deviation gains.}
Fix a persistent agent identity, information and feasible-strategy spaces, and an initial history preceding
strategic choices. Write the added reward as \(\rho F_t\), where
\[
F_t=\gamma\Phi_{t+1}-\Phi_t.
\]
For a horizon \(H\), shifting indices gives
\[
\sum_{t=0}^{H-1}\gamma^t F_t
=-\Phi_0+\gamma^H\Phi_H=-\Phi_0
\qquad (\Phi_H=0).
\]
If \(\Phi_0\) is fixed under deviations, every strategy's payoff changes by the same constant. Thus no
profitable joint report/action deviation becomes unprofitable: exact potential-based shaping cannot newly
implement Q's truthfulness or obedience requirements. This permits learning improvements without new strategic
incentives (ledger \#1, \#10, \#14; ada's boundary note).

\paragraph{Participation masks can leave transfers.}
Let \(m_t\in\{0,1\}\) indicate whether the agent's return includes the shaping reward at transition \(t\). The same index shift now yields
\[
\begin{aligned}
\sum_{t=0}^{H-1}\gamma^t m_t F_t
={}&-m_0\Phi_0
+\sum_{t=1}^{H-1}\gamma^t(m_{t-1}-m_t)\Phi_t\\
&+\gamma^H m_{H-1}\Phi_H.
\end{aligned}
\]
With zero terminal potential and zero potential while inactive, exit terms vanish, but each re-entry at time
\(e\) leaves a debit \(-\gamma^e\Phi_e\). Coalition changes alone do not cause this problem; omitted
increments do (ledger \#8, \#13--14).

Ada's exact example takes \(\gamma=0.99\), \(\rho=1\), zero base rewards, mask \((1,0,1)\), and potentials
\((1/2,0,p,0)\). Its masked return is
\[
G_{\mathrm{mask}}(p)=-\frac12-(0.99)^2p.
\]
An explicitly assumed evidence-withholding action lowers the returning score from \(p=1/2\) to \(p=1/10\)
without changing the physical trajectory. The resulting gain is
\[
G_{\mathrm{mask}}(1/10)-G_{\mathrm{mask}}(1/2)
=(0.99)^2(0.4)=0.39204.
\]
Full accounting gives \(-1/2\) for either choice. Under zero inactive potentials, adding \(\Phi_e\) at
re-entry cancels the debit; generally the interior settlement is \((m_t-m_{t-1})\Phi_t\). Ada's script checks
11,664 rational instances of the identity and settlement, and emmy independently verified the checks and
example (ledger \#13--14; ada's note). This is a conditional counterexample to masked invariance, not an
observed exploit in P.

\paragraph{Repeated judgments cannot replace missing evidence.}
Emmy fixes a type, report, recommendation, and other agents' strategies. Let a feasible deviation have
intrinsic expected gain \(g>0\); let \(Z\) contain all judge evidence, with desired/deviating laws
\(\mu_\star,\mu_d\). A transcript \(Y_K\) generated by a common stochastic kernel from this fixed evidence has
laws \(\nu_\star^K,\nu_d^K\). For an additive transfer \(T(Y_K)\) of range at most \(B\), with no other
payment or prohibition channel, obedience requires
\[
\begin{aligned}
g&\leq E_{\nu_\star^K}[T]-E_{\nu_d^K}[T]\\
 &\leq B\operatorname{TV}(\nu_\star^K,\nu_d^K)
 \leq B\operatorname{TV}(\mu_\star,\mu_d).
\end{aligned}
\]
Subtracting the transfer's infimum and dividing by \(B\) gives a function in \([0,1]\), whose expectation
difference is bounded by total variation. Conditioning that function on \(Z\) proves the last inequality.
Identical evidence laws therefore leave a profitable deviation undeterred, regardless of query count. For a
single deviation, paying \(B\) where the desired transcript law exceeds the deviation law attains the
transcript bound; multiple deviations need a common contract (ledger \#6, \#12, \#16; emmy's information
note).

The synthetic example uses desired/deviating probabilities \(0.6/0.4\) for binary evidence \(Z=1\),
conditional judgment probabilities \(0.9/0.1\) for \(Z=1/0\), \(g=0.3\), and \(B=1\). Reusing one image gives
\[
\operatorname{TV}(\nu_\star^K,\nu_d^K)
=0.2\operatorname{TV}\bigl(\operatorname{Bin}(K,0.9),
\operatorname{Bin}(K,0.1)\bigr)\leq0.2.
\]
Thus no query count suffices. Fresh independent evidence instead gives judgment probabilities \(0.58/0.42\).
At \(K=8\), shared-image TV is \(0.1989088\), whereas fresh-image TV is \(0.3411767867392>g/B\), sufficient
for the single-deviation likelihood-ratio payment. Both peers checked these finite-distribution calculations.
They illustrate a possible non-shaping contract, not P's objective or benchmark performance (ledger \#12,
\#16).

\section*{Objections and their disposition}
The invariance statement needed qualifications: rewards must not enlarge the information or strategy spaces,
the initial potential must precede strategic choice, and adjacent increments must reuse the same realized
potential. Otherwise an initial report can leave a transfer \(-\rho E[\Phi_0(\hat t)]\), or independently
resampled potentials can leave residuals. These qualifications were accepted; they do not establish a
reporting channel in P (ledger \#10, \#14).

Ada's suggestion that common-score translation destroys implementation was narrowed after emmy's objection.
Bradley--Terry probabilities and softmax do lose that distinction, but this matters only if indistinguishable
cases require different incentives or contain a profitable hidden deviation. Q itself can use identifying
action differences. Ada accepted this objection, so translation invariance is not a standalone impossibility
result (ledger \#5, \#11, \#16).

The evidence ceiling assumes repeated queries acquire no new evidence; fresh views change that premise (ledger
\#10, \#12). The participation attack assumes a feasible, sufficiently inexpensive manipulation of evidence
and a particular reward mask. P's inactive-agent accounting and singleton conventions remain unspecified in
the supplied material, so the empirical objection is unresolved (ledger \#8, \#13--14).

\section*{Sources outside Q and P}
The peers used Devlin and Kudenko (2012), on dynamic potential-based shaping preserving Nash equilibria, and
Grzes (2017), on episodic terminal-potential effects and their correction (ledger \#7, \#13). They also used
Saig, Einav, and Talgam-Cohen, \emph{Incentivizing Quality Text Generation via Statistical Contracts}, NeurIPS
2024, Proposition~1: its characterization using TV distance to mixtures of deviations subsumes the two-action
budget bound (ledger \#9). The peers' notes record the source URLs and searches. These checks remove novelty
claims for the generic invariance and TV arguments; they do not settle prior work on the particular
participation audit (ledger \#13--14, \#16).

\section*{What remains open}
The material is thin on implementation and empirical behavior. No runnable P implementation was included or
audited. Actual respawn identity, reward accounting, stored potentials, terminal handling, feasible view
manipulation, and strategic utilities remain unresolved. P uses comparisons during training; a strategic
deployment mechanism would be an extension. The evidence supports a specified investigation, not a completed
publication-level result or a finding that P lacks its claimed learning benefits (ledger \#14, \#16; both
peers' notes).

\section*{Smallest next step}
Inspect and instrument one disappearance/re-entry trace in a runnable P implementation, recording persistent
identity, potentials, included shaping increments, and the discounted return. Compare its accumulated shaping
with the boundary identity and full accounting. This settles whether the proposed accounting residual occurs
in that implementation. Only if it does should an explicitly feasible evidence intervention test whether it
creates a profitable deviation; correct accounting or an unavailable intervention defeats the corresponding
attack prediction. That staged stopping criterion is the next step specified in ada's note and endorsed by
emmy (ledger \#13--14, \#16).

\end{document}
